% 图契约：本题的源、边界和推导逐项对应；箭头符号按图注解释。
\documentclass[tikz,border=7pt]{standalone}
\usepackage{fontspec,xeCJK,amsmath,amssymb}
\IfFontExistsTF{Times New Roman}{\setmainfont{Times New Roman}}{\setmainfont{TeX Gyre Termes}}
\IfFontExistsTF{Heiti SC}{\setCJKmainfont[AutoFakeBold=1.5]{Heiti SC}}{\setCJKmainfont[AutoFakeBold=1.5]{FandolHei-Regular}}
\usetikzlibrary{arrows.meta,bending,calc,patterns,positioning}
\definecolor{ink}{HTML}{243745}
\definecolor{blue}{HTML}{275DAB}
\definecolor{green}{HTML}{177C63}
\definecolor{red}{HTML}{BD4935}
\definecolor{orange}{HTML}{AA6C16}
\definecolor{purple}{HTML}{715393}
\tikzset{every node/.style={font=\fontsize{9.3}{12.5}\selectfont,text=ink},
axis/.style={-{Stealth[length=4pt]},line width=.65pt,draw=ink!65},
vec/.style={-{Stealth[length=5pt]},line width=1pt,draw=blue},
short/.style={-{Stealth[length=3pt]},line width=.8pt},
guide/.style={dashed,draw=ink!40,line width=.5pt},
card/.style={rounded corners=2pt,draw=blue!65,fill=blue!4,align=center,text width=6.5cm,inner xsep=9pt,inner ysep=7pt},
title/.style={font=\fontsize{10.5}{14}\selectfont\bfseries,text width=15cm,align=center},
label/.style={font=\fontsize{8.4}{11}\selectfont,align=center},
note/.style={label,text width=6.2cm}}
\begin{document}
\begin{tikzpicture}[x=1cm,y=1cm]
\path[use as bounding box] (-8,-3.5) rectangle(8,3.5);
\node[title] at (0,3.15) {M2-140\quad 半球和平面要同时成为零势面};
\begin{scope}[shift={(-4,0)}]

\fill[ink!6](-2.4,-2)rectangle(2.4,0);\fill[ink!6](0,0)circle(1);\draw[ink,thick](-2.4,0)--(-1,0) (1,0)--(2.4,0);\draw[ink,thick](1,0)arc(0:180:1);
\fill[red](0,1.9)circle(2pt)node[right]{$Q$};\fill[blue](0,-1.9)circle(2pt)node[right]{$q_1$};\fill[blue](0,.526)circle(2pt)node[right]{$q_2$};\fill[red](0,-.526)circle(2pt)node[right]{$q_3$};
\node[note] at (0,-2.85) {灰区不属于真实求解域，三枚像源全部位于其中。};
\end{scope}
\node[card,draw=blue!65,fill=blue!4] at (3.9,1.8) {$q_1=-Q,\ z_1=-d$};
\node[card,draw=green!65,fill=green!4] at (3.9,0) {$q_2=-QR/d,\ z_2=R^2/d$\\$q_3=+QR/d,\ z_3=-R^2/d$};
\node[card,draw=orange!65,fill=orange!4] at (3.9,-1.8) {反射负责平面，反演负责半球；\\只保留一枚像源无法同时满足。};
\end{tikzpicture}\end{document}